\section{评估系统设计：从单点评分到 Holistic Evaluation}

训练数据只有通过稳定测量才能形成迭代闭环。本节承接官方 slides 13--18，把“评估一个模型”拆成任务、harness、工具空间、指令鲁棒性和 benchmark 质量五个层面，重点说明模型能力与脚手架能力如何分离，以及单一平均分为何不足以定义智能。

\subsection{Verifier 不是一个分数，而是一组机制}
上一节从数据角度定义 verifier，这里转向评估系统：同一轨迹需要结果、结构、过程和安全多层检查。只有把这些检查的触发条件和失败原因分别记录，团队才能知道模型是不会解题、不会用工具，还是利用了测试漏洞。
课程中一个很重要的观点是：verifier 不应被理解为单个布尔规则，而应是多维度验证体系。代码任务可能用单元测试，数学任务可用 proof checker，结构化输出任务要校验 JSON 合法性与字段一致性，业务任务还要检查副作用边界。

\begin{knowledgebox}{Verifier 的分层设计}
\begin{itemize}
    \item L1 结果层：最终答案是否正确。
    \item L2 结构层：输出是否满足协议与格式要求。
    \item L3 过程层：关键中间步骤是否符合策略约束。
    \item L4 安全层：是否触发越权调用、数据泄露或危险动作。
\end{itemize}
\end{knowledgebox}

\subsection{Harness 交换与鲁棒性验证}
讲者提到应在评估中引入 harness swapping：同一任务在不同执行框架、不同工具包装、不同系统提示下重复评测。如果性能大幅波动，往往说明模型依赖了偶然实现细节而非稳健策略。

\begin{importantbox}{Harness 交换的价值}
它不是 ``再测一次''，而是检测模型是否具备环境迁移鲁棒性。若模型只在某个固定 harness 上表现优秀，部署风险会很高。
\end{importantbox}

\subsection{评估污染与难度监控}
课程中还强调 benchmark contamination 与 hardness drift。随着社区迭代加速，公开数据可能被训练过程吸收，导致指标失真。评估系统必须持续监控题目难度与分布变化。

\begin{warningbox}{高分不等于高能力}
当 benchmark 被污染或难度下降时，高分可能只意味着 ``见过类似题''。因此评估要包含：
\begin{itemize}
    \item 新样本引入频率控制；
    \item 题目泄露风险审计；
    \item 难度分层报告（easy/medium/hard）；
    \item 跨 benchmark 交叉验证。
\end{itemize}
\end{warningbox}

\subsection{本章小结}
可验证 Agent 的评估必须是 Holistic 的：多 verifier、多 harness、多难度分层、持续污染监控。单一分数无法支撑可靠部署。

